A magnitude do torque no aparato do Problema 3.20 depende da configuração geométrica e das propriedades gravitacionais.

\textbf{Caso de massa pontual} a distância $R$ do eixo:
\[
\tau = m\,g\,R,
\]
onde $m$ é a massa aplicada, $g$ a aceleração gravitacional e $R$ o braço de momento.

\textbf{Caso de roda/disco com massa distribuída:}
Se a roda tem distribuição de massa com elementos $dm_i$ localizados a braços perpendiculares $r_{\perp,i}$ em relação ao eixo horizontal:
\[
\tau = \sum_i g\,r_{\perp,i}\,dm_i = g\int r_\perp\,dm.
\]
Para uma roda de massa total $M_w$ cujo centro de massa está a uma distância $d_{\rm cm}$ do eixo:
\[
\tau = M_w\,g\,d_{\rm cm}.
\]

\textbf{Relação com as propriedades do aparato:}
O torque gerado é proporcional ao \textbf{peso} da configuração ($M_w\,g$) e ao \textbf{braço geométrico} ($R$ ou $d_{\rm cm}$), confirmando que a magnitude do torque depende exclusivamente de grandezas gravitacionais e geométricas fixas do aparato.

